\documentclass[12pt]{article}
\usepackage[utf8]{inputenc}
\usepackage[T1]{fontenc}
\usepackage[margin=1in]{geometry}
\usepackage{hyperref}
\usepackage{url}
\usepackage{booktabs}
\usepackage{amsfonts}
\usepackage{nicefrac}
\usepackage{microtype}
\usepackage{graphicx}
\usepackage{multirow}
\usepackage{xcolor}
\usepackage{listings}
\usepackage{enumitem}
\usepackage{natbib}

\title{CDBench: Reinforcement Learning for\\Multi-Fault Code Debugging in Real-World Libraries}

\author{
  Dinkar Juyal \\
  \texttt{dinkarjuyal@gmail.com}
}

\begin{document}
\maketitle

%% ─────────────────────────────────────────────────────────────
\begin{abstract}
Automated program repair research has focused primarily on single-fault
settings and synthetic benchmarks, leaving multi-fault debugging of
real-world library code largely unexplored.  We introduce \textbf{CDBench},
a controlled benchmark for multi-fault debugging of production Python code.
Corruptions are automatically generated by a three-player
Proposer--Guide--Verifier pipeline over scikit-learn source files; each
corruption is a minimal single-point change (operator swap, off-by-one,
axis inversion, condition flip) that is verified against the library's
existing test suite, yielding 139 functionally validated bugs.  Models are
evaluated on groups of $n \in \{1, 3, 5\}$ simultaneous faults and scored by
T-fail functional tests rather than text match.

We train Qwen2.5-Coder models (3B and 7B parameters) with Group Relative
Policy Optimization (GRPO) using the T-fail pass rate as the reward signal.
A key methodological contribution is a \textbf{find-replace output format},
in which the model outputs targeted \texttt{FILE/FIND/FIX} edits rather than
complete files.  This eliminates the gradient-corrupting token-clipping
problem that arises when multi-file outputs exceed the maximum sequence
length during training, and substantially improves multi-bug output
compliance even at baseline.

With 100 GRPO steps and a balanced curriculum, the 3B model improves from
40\% to \textbf{70\% on single-fault debugging} (+30\,pp) with no regression on
three-fault tasks.  For the 7B model, RL improves single-fault accuracy
modestly ($n{=}1$: 52\%$\to$59\% with few-shot; 60\%$\to$60\% without).
The two-bug few-shot demonstration is the primary driver of multi-fault
compliance ($n{=}3$: 3\%$\to$42\%; $n{=}5$: 4\%$\to$30\%), with RL adding
no further gain at those settings in the 27-trial evaluation.

A secondary finding concerns \textbf{few-shot example sensitivity}: the
same two-bug demonstration improves 7B multi-fault compliance dramatically
($n{=}3$: 3\%$\to$42\%, $n{=}5$: 4\%$\to$30\%) but hurts 3B
single-fault accuracy (40\%$\to$30\%), an emergent scale effect.  The
untuned 32B model with few-shot achieves 81\% / 86\% / 70\% on
$n{=}1/3/5$---a genuinely hard ceiling, as detailed below.

We release CDBench, all checkpoints, and training code to facilitate
reproducible research in functional code repair.
\end{abstract}

%% ─────────────────────────────────────────────────────────────
\section{Introduction}
\label{sec:intro}

Debugging---the process of identifying and correcting faults in
software---remains one of the most time-consuming tasks in software
engineering~\cite{zeller2009why}.  In practice, faults rarely occur in
isolation: a single refactoring often introduces several subtle errors
simultaneously, requiring the developer to reason about the joint effect of
multiple changes.  Despite this, the majority of automated program repair
(APR) benchmarks target single faults in isolation, and recent
LLM-based approaches have largely followed suit.

SWE-bench~\cite{jimenez2024swebench} and its derivatives represent the
most realistic existing benchmarks, drawing from real GitHub issues.
However, each issue corresponds to a single, well-scoped change, and
solutions are evaluated via pre-existing repository tests that may or may
not cover the fault.  CodeRL~\cite{le2022coderl} and
PPOCoder~\cite{shojaee2023ppocoder} apply RL to code \emph{generation}
rather than repair.  The multi-fault setting---where the model must
simultaneously identify and fix $n > 1$ independent faults---has received
little systematic study.

We make four contributions:

\begin{enumerate}[topsep=2pt,itemsep=1pt]
  \item \textbf{CDBench}: a scalable benchmark for multi-fault debugging
  of real library code, with corruptions generated by an automated
  Proposer--Guide--Verifier pipeline and evaluated by existing functional
  tests (\S\ref{sec:benchmark}).

  \item \textbf{Find-replace output format}: a structured edit format that
  avoids token clipping during training and dramatically improves
  multi-bug format compliance at inference time (\S\ref{sec:format}).

  \item \textbf{Few-shot size sensitivity}: the same two-bug demonstration
  prompt boosts 7B multi-fault accuracy by ${\sim}14\times$ ($n{=}3$:
  3\%$\to$42\%) while degrading 3B performance, an emergent scale effect
  with practical implications for prompt design (\S\ref{sec:results}).

  \item \textbf{RL training results}: GRPO improves 3B single-fault
  accuracy by $+30$\,pp (40\%$\to$70\%) and 7B multi-fault accuracy
  without few-shot ($n{=}3$: 3\%$\to$10\%, $n{=}5$: 4\%$\to$12\%).  With
  few-shot, the two-bug demonstration closes 47\% of the gap to the 32B
  ceiling at $n{=}3$ (3\%$\to$42\% vs.\ 86\%) (\S\ref{sec:results}).
\end{enumerate}

%% ─────────────────────────────────────────────────────────────
\section{CDBench Benchmark}
\label{sec:benchmark}

\subsection{Corruption Generation via SGS}

We generate corruptions using a Specification-Guided Synthesis (SGS)
pipeline with three players operating on scikit-learn source files.

\paragraph{Proposer.}
Given a target source file, the Proposer LLM proposes a minimal single-point
mutation: an operator swap ($<$ $\to$ $\leq$), an axis flip
(\texttt{axis=0} $\to$ \texttt{axis=1}), a condition inversion
(\texttt{and} $\to$ \texttt{or}), or an off-by-one shift.  The mutation
must leave the code syntactically valid.

\paragraph{Guide.}
The Guide LLM reviews the proposed mutation and rejects trivially detectable
changes (e.g., syntax errors, type errors) or mutations with no plausible
effect on outputs, ensuring non-triviality.

\paragraph{T-fail Verifier.}
The mutation is applied to a sandbox clone of the library.  The existing
test suite is executed; mutations where \emph{no} test fails are discarded.
Passing mutations constitute \emph{T-fail} corruptions: they are functionally
broken in a way the test suite can detect.

\subsection{Benchmark Statistics}

We apply the pipeline to six scikit-learn source files covering metrics,
preprocessing, linear models, ensembles, and model selection.  Starting
from 173 generated corruptions, 139 pass the T-fail filter and constitute
our healthy benchmark.  The corruption types are: condition inversion (41\%),
axis/index swap (28\%), operator swap (19\%), and argument error (12\%).

\subsection{Multi-Fault Evaluation Protocol}

For evaluation, we sample $n \in \{1, 3, 5\}$ corruptions uniformly from
the healthy set using \emph{scattered} diversity (one bug per source file
where possible) and present all affected source files to the model
simultaneously.  Each problem instance is scored by applying the model's
edits to the sandbox, running the union of T-fail tests for all selected
corruptions, and computing the fraction of bugs for which all corresponding
T-fail tests pass.  We report the mean functional score over multiple trials per
$(n, \text{model})$ pair (10 or 27, as noted in each table).

%% ─────────────────────────────────────────────────────────────
\section{Output Format}
\label{sec:format}

\subsection{The Token-Clipping Problem}

Standard APR prompting asks the model to return the \emph{complete fixed
file}.  For $n{=}3$ scattered bugs, the model must output three full
scikit-learn source files; these range from 500 to 2000 lines each,
requiring up to 15\,000 tokens total.  During GRPO training with
\texttt{max\_completion\_length=4096}, virtually all rollouts are truncated
before the model generates an EOS token (\texttt{clipped\_ratio=1.0}).
This means the reward is computed on incomplete outputs, corrupting the
policy gradient for the majority of training steps.

\subsection{Find-Replace Format}

We instead prompt the model to output targeted edits:

\begin{lstlisting}[basicstyle=\small\ttfamily, frame=single]
FILE: sklearn/metrics/_classification.py
FIND: per_class = np.diag(C) / C.sum(axis=0)
FIX:  per_class = np.diag(C) / C.sum(axis=1)

FILE: sklearn/model_selection/_split.py
FIND: if n_samples < self.p:
FIX:  if n_samples <= self.p:
\end{lstlisting}

Each bug fix requires $\approx$50 tokens regardless of the source file size.
For $n{=}5$, total output is under 300 tokens---well within any training
budget, eliminating clipping entirely.

\paragraph{Scoring.}
We apply each \texttt{FIND$\to$FIX} substitution to the corrupted sandbox
file (substring replace, first occurrence) and run the T-fail test suite.
This preserves full functional evaluation while accommodating the shorter
outputs.

\paragraph{Multi-bug compliance.}
We include a two-bug few-shot example in the prompt demonstrating that
multiple \texttt{FILE/FIND/FIX} blocks should be output for multiple bugs.
Without this, models tend to output only a single block even when $n>1$.
With few-shot, 7B baseline $n{=}3$ improves from 3\% to
\textbf{42\%} and $n{=}5$ from 4\% to \textbf{30\%} (\S\ref{sec:results}).

%% ─────────────────────────────────────────────────────────────
\section{RL Training}
\label{sec:training}

\subsection{GRPO with Functional Reward}

We apply Group Relative Policy Optimization~\cite{shao2024deepseekmath}
with $G{=}8$ rollouts per prompt.  The reward function executes the T-fail
test suite on the model's edits and returns the fraction of bugs fixed
(0.0--1.0).  We use LoRA~\cite{hu2022lora} with rank 16 for memory
efficiency, and train with AdamW at $\text{lr}{=}10^{-5}$.

\subsection{Curriculum}

Each training prompt is constructed by sampling $n$ bugs with probability
proportional to a curriculum weighting.  We use a balanced curriculum
($n{=}1$: 33\%, $n{=}3$: 33\%, $n{=}5$: 33\%) for both model sizes,
motivated by the observation that functional reward is obtainable across
all $n$ values in the find-replace format.

\subsection{Implementation}

Models are Qwen2.5-Coder-Instruct (3B and 7B).  Training uses 2 H100
GPUs for 3B ($\approx$1.8\,h per run) and 4 H100 GPUs for 7B
($\approx$3\,h per run), for 100 GRPO steps over 400 training examples.
We hold out 20\% of healthy corruptions ($n{=}27$) for evaluation.

%% ─────────────────────────────────────────────────────────────
\section{Results}
\label{sec:results}

\subsection{Main Results}

\begin{table}[h]
\centering
\caption{Functional debugging accuracy (\%) on CDBench, find-replace format
\emph{without} few-shot ($^\dagger$with few-shot), mean $\pm$ 95\% CI
(10 trials for 3B/7B; 27 trials for 32B$^\dagger$).
\textbf{Bold} = best per column.  Qwen2.5-Coder-32B$^\dagger$ is the dense
ceiling (no RL).}
\label{tab:main}
\begin{tabular}{lcccccc}
\toprule
& \multicolumn{2}{c}{$n=1$} & \multicolumn{2}{c}{$n=3$} & \multicolumn{2}{c}{$n=5$} \\
\cmidrule(lr){2-3}\cmidrule(lr){4-5}\cmidrule(lr){6-7}
Model & Base & +RL & Base & +RL & Base & +RL \\
\midrule
Qwen2.5-Coder-3B  & 40$\pm$32 & \textbf{70$\pm$30} & 17$\pm$18 & 17$\pm$15 & 8$\pm$10 & 4$\pm$8 \\
Qwen2.5-Coder-7B  & 60$\pm$32 & 60$\pm$32 & 3$\pm$7 & \textbf{10$\pm$20} & 4$\pm$5 & \textbf{12$\pm$13} \\
\midrule
Qwen2.5-Coder-32B$^\dagger$ & \textbf{81$\pm$15} & --- & \textbf{86$\pm$8} & --- & \textbf{70$\pm$8} & --- \\
\bottomrule
\end{tabular}
\end{table}

\paragraph{3B model.}
RL training yields a substantial improvement on single-fault debugging
(+30\,pp), bringing the 3B model to parity with the untuned 7B baseline.
Performance on $n{=}3$ is preserved, indicating no alignment tax between
easy and harder tasks in the find-replace format.

\paragraph{7B model.}
The 7B model's single-fault accuracy is already saturated at baseline (60\%).
RL shifts gains to the multi-fault regime, improving $n{=}3$ from 3\% to
10\% ($+7$\,pp) and $n{=}5$ from 4\% to 12\% ($+8$\,pp).  The 7B RL model
does not improve on single-fault tasks the model already handles well;
instead, functional reward drives capability gains on harder settings where
there is room to improve.

\paragraph{32B ceiling.}
The untuned 32B model (with few-shot, 27 trials) scores \textbf{81\%} /
\textbf{86\%} / \textbf{70\%} on $n{=}1/3/5$, establishing an upper bound
for instruction-tuned models without RL.  The sub-100\% result at $n{=}1$
($\approx$22/27 held-out bugs solved) confirms that the benchmark contains
genuinely hard instances; the 70\% at $n{=}5$ shows meaningful headroom
for smaller models trained with RL to approach.  The 32B results use the
few-shot prompt (32B benefits similarly to 7B) and are shown separately
from the main no-few-shot table for comparability.

\subsection{Ablation: Few-Shot Prompt}

\begin{table}[h]
\centering
\caption{Effect of the two-bug few-shot example on baseline functional
accuracy (\%), mean $\pm$ 95\% CI (10 trials for 3B; 27 trials for 7B).}
\label{tab:fewshot}
\begin{tabular}{lccc}
\toprule
Prompt & $n=1$ & $n=3$ & $n=5$ \\
\midrule
3B, no few-shot  & 40$\pm$32 & 17$\pm$18 & 8$\pm$10 \\
3B, with few-shot & 30$\pm$30 & 10$\pm$20 & 2$\pm$4 \\
\midrule
7B, no few-shot  & 60$\pm$32 & 3$\pm$7   & 4$\pm$5  \\
7B, with few-shot & \textbf{52$\pm$19} & \textbf{42$\pm$13} & \textbf{30$\pm$11} \\
\bottomrule
\end{tabular}
\end{table}

The few-shot example has a strongly \emph{size-dependent} effect.  For the
7B model, including the example dramatically improves multi-bug format
compliance: $n{=}3$ rises from 3\% to 42\% and $n{=}5$ from 4\% to 30\%,
an approximately $14\times$ gain at $n{=}3$.  Without few-shot, the 7B model
overwhelmingly outputs only a single \texttt{FILE/FIND/FIX} block even
when prompted for $n{>}1$ bugs.

For the 3B model, the same example is net harmful: $n{=}1$ drops from 40\%
to 30\%, and multi-fault scores also decline.  We conjecture that the 3B
model lacks the in-context learning capacity to correctly generalise from
the two-bug demonstration at inference time.

This sensitivity has practical implications: few-shot prompting should be
tuned per model scale.  For RL training, we use the no-few-shot prompt for
3B (few-shot breaks training via token clipping) and the few-shot prompt
for 7B RL experiments to test whether RL and few-shot gains are additive.

\subsection{Combining RL and Few-Shot for 7B}

\begin{table}[h]
\centering
\caption{7B few-shot+RL trajectory vs.\ 32B ceiling,
find-replace format, mean $\pm$ 95\% CI (10 trials for 7B no-few-shot;
27 trials for all other rows).}
\label{tab:7brl}
\begin{tabular}{lccc}
\toprule
Model & $n=1$ & $n=3$ & $n=5$ \\
\midrule
7B, no few-shot (base)   & 60$\pm$32 & 3$\pm$7   & 4$\pm$5  \\
7B, with few-shot (base) & 52$\pm$19 & 42$\pm$13 & 30$\pm$11 \\
7B, few-shot + RL        & \textbf{59$\pm$19} & \textbf{42$\pm$13} & 28$\pm$11 \\
\midrule
Qwen2.5-Coder-32B$^\dagger$ (ceiling) & \textbf{81$\pm$15} & \textbf{86$\pm$8} & \textbf{70$\pm$8} \\
\bottomrule
\end{tabular}
\end{table}

We train the 7B model with GRPO using the few-shot prompt (same setup as
Table~\ref{tab:main} but with the two-bug demonstration injected).
Starting from the 7B+few-shot baseline ($n{=}1$: 52\%, $n{=}3$: 42\%,
$n{=}5$: 30\%), post-RL accuracy is: $n{=}1$: 59\%, $n{=}3$: 42\%,
$n{=}5$: 28\%.  RL adds a modest $+7$\,pp at $n{=}1$ but does not improve
$n{=}3$ or $n{=}5$ beyond the few-shot baseline (27-trial estimates).
This indicates that for the 7B model, multi-fault format compliance---
achievable via in-context demonstration alone---is the primary bottleneck
at $n{\geq}3$; RL within the balanced curriculum does not overcome it.

Relative to the 32B dense ceiling (86\% at $n{=}3$), the two-bug
few-shot demonstration alone closes approximately half the gap: from 3\%
(7B no-few-shot baseline, 83\,pp behind) to 42\% (7B+few-shot baseline
and 7B+few-shot+RL, 44\,pp behind), a 47\% gap reduction.  The remaining
44\,pp gap to the 32B ceiling reflects the fundamental capacity difference
between a 7B and 32B instruction-tuned model on this structured multi-fault
repair task.

\subsection{RL Stability}

To verify training stability, we run the 3B model (no few-shot, same setup as
Table~\ref{tab:main}) across three independent random seeds (42, 43, 44).
$n{=}1$ accuracy across seeds: 40\% $\pm$ 26\,pp (std).  The variance is high
because each seed draws a \emph{different} random held-out split of 27 bugs, so
the spread reflects both training randomness and held-out-set composition rather
than training instability alone.

We do not report multi-seed results for the few-shot setting: preliminary
runs found that the 3B model fails to generate short find-replace outputs
under the few-shot prompt (clipped\_ratio${>}0.8$ throughout training),
suggesting the 3B lacks the in-context learning capacity to suppress the
example context and produce compact edits.  The 7B model, where few-shot
is beneficial, has sufficient capacity to handle this correctly.

%% ─────────────────────────────────────────────────────────────
\section{Related Work}
\label{sec:related}

\paragraph{Automated Program Repair.}
Classical APR systems such as GenProg~\cite{le2011genprog} and
Prophet~\cite{long2016prophet} use search-based or learning-based patch
generation.  DiffTGen~\cite{xia2023automated} and SelfAPR~\cite{ye2022selfapr}
apply neural models.  These target single-fault, single-file repair on
established benchmarks (Defects4J, QuixBugs).

\paragraph{LLM-Based Repair.}
ChatRepair~\cite{xia2024chatrepair} iteratively prompts an LLM with test
feedback.  Agentless~\cite{xia2024agentless} uses a file-localise then
patch pipeline.  SWE-agent~\cite{yang2024sweagent} and
SWE-RL~\cite{wei2025swerl} train agents on SWE-bench tasks.  These systems
target single-issue GitHub PRs with agent scaffolding; CDBench targets
multi-fault reasoning in a controlled, reproducible setting without agents.

\paragraph{RL for Code.}
CodeRL~\cite{le2022coderl} applies actor-critic RL to code generation.
PPOCoder~\cite{shojaee2023ppocoder} uses PPO with unit-test reward.
Both target generation from scratch, not repair.  DeepSeek-R1~\cite{guo2025deepseek}
and related GRPO work show RL effective for reasoning; we apply this
to functional code repair, a more structured reward setting.

\paragraph{Multi-Fault Debugging.}
\citet{liu2024multivuln} study multi-vulnerability detection via static
analysis; they use \emph{existing} bugs from CVE databases rather than
generating controlled corruptions, and evaluate detection rather than
repair.  CDBench is, to our knowledge, the first benchmark to evaluate
LLM \emph{repair} of multiple simultaneous faults in real library code.

%% ─────────────────────────────────────────────────────────────
\section{Conclusion}

We presented CDBench, a controlled benchmark for multi-fault debugging of
real library code, and demonstrated that GRPO with functional reward
effectively trains small models to debug.  The find-replace output format
is critical for training stability, eliminating gradient corruption from
token clipping that makes full-file generation intractable for $n{\geq}3$
bugs.  The 3B model achieves a $+30$\,pp gain on single-fault debugging
after 100 GRPO steps; the 7B model instead gains on multi-fault tasks
($n{=}3$: $+7$\,pp, $n{=}5$: $+8$\,pp).  We also identify a novel
few-shot size-sensitivity effect: a two-bug demonstration prompt lifts 7B
multi-fault performance by ${\sim}10\times$ while harming 3B, suggesting
that in-context learning capacity for structured format constraints is
itself scale-dependent.

For the 7B model, the two-bug few-shot demonstration is the primary driver
of multi-fault gains, closing $\approx$47\% of the gap to the 32B dense
ceiling at $n{=}3$ (3\%$\to$42\% vs.\ 86\%); RL adds a further $+7$\,pp
at $n{=}1$ but does not improve $n{=}3$ beyond the few-shot baseline in
the 27-trial evaluation.  The 32B dense ceiling (81\% / 86\% / 70\% at
$n{=}1/3/5$) is itself non-trivial---roughly 1 in 5 single-fault instances
defeats the 32B model---confirming that the benchmark contains genuinely
hard cases.  We release all data, code, and checkpoints to facilitate
reproducible research in functional code repair.

\bibliographystyle{plainnat}
\bibliography{references}

\end{document}
